\documentclass[a4paper]{article}
\usepackage{amsfonts}
\usepackage{amsmath}
\usepackage{amssymb}
\usepackage{graphicx}     

\begin{document}
	
\noindent \large Auteur: Dara van Engelen     \\
\noindent \large Studierichting: Informatica \\
\noindent \large Oefeningenreeks 6B nr. 5.6 \\
	
\medskip
	
\normalsize

Bereken $q$ als $x_1 = 3$ en $s_4 = -15$\\

$s_n = x_1 \cdot \dfrac{1-q^n}{1-q}$\\

$\Rightarrow -15 = 3 \cdot \dfrac{1-q^4}{1-q} \Leftrightarrow -5 = \dfrac{1-q^4}{1-q} \Leftrightarrow -5 + 5q = 1 - q^4$\\

$\Leftrightarrow q^4 + 5q -6 = 0 \Leftrightarrow (q-1)(q+2)(q^2-q+3) = 0$\\

$q = 1$ of $ q = -2$\\

$ q = 1$: constante rij dus $s_4 = 4\cdot3 = 12 \neq -15$ dus kan niet\\

$ q = -2$ is de oplossing.


\begin{thebibliography}{999}
%\bibitem{a} Referentie
\end{thebibliography}
\end{document}